\documentclass[UTF8,a4paper,11pt]{ctexart}
\usepackage[margin=22mm,headheight=14pt,footskip=12mm]{geometry}
\usepackage{amsmath,amssymb,array,longtable,booktabs,tikz,needspace,fancyhdr,hyperref}
\usepackage{microtype}
\IfFontExistsTF{MiSans}{\setCJKmainfont[BoldFont=MiSans Bold]{MiSans}\setCJKsansfont[BoldFont=MiSans Bold]{MiSans}}{}
\hypersetup{colorlinks=false,pdfborder={0 0 0},pdfauthor={Good Learning}}
\newcolumntype{P}[1]{>{\raggedright\arraybackslash}p{#1}}
\setlength{\LTleft}{0pt}\setlength{\LTright}{0pt}
\setlength{\parindent}{0pt}\setlength{\parskip}{.65em}
\linespread{1.18}
\raggedbottom
\setcounter{secnumdepth}{0}
\ctexset{section={format=\Large\bfseries\sffamily,beforeskip=1.35em,afterskip=.65em},
subsection={format=\large\bfseries\sffamily,beforeskip=1em,afterskip=.45em}}
\newcommand{\chaptertitle}[1]{\Needspace{8\baselineskip}\noindent{\LARGE\bfseries\sffamily #1}\par\vspace{4mm}\hrule\vspace{6mm}}
\newcommand{\CompactQuestion}[2]{\par\noindent\textbf{#1.}\ #2\par\vspace{.8em}}
\newcommand{\PracticeQuestion}[3]{\par\Needspace{4\baselineskip}\noindent\begin{minipage}{\linewidth}\textbf{#1.}\ #2\par\noindent\rule{0pt}{#3mm}\par\end{minipage}\par}
\newcommand{\CompactQuestionColon}[2]{\par\noindent\textbf{#1：}\ #2\par\vspace{.8em}}
\newcommand{\PracticeQuestionColon}[3]{\par\Needspace{4\baselineskip}\noindent\begin{minipage}{\linewidth}\textbf{#1：}\ #2\par\noindent\rule{0pt}{#3mm}\par\end{minipage}\par}
\pagestyle{fancy}\fancyhf{}\fancyfoot[C]{\thepage}\renewcommand{\headrulewidth}{0pt}
\begin{document}
\thispagestyle{empty}
\noindent{\sffamily\footnotesize 每日课程}\hfill{\sffamily\footnotesize 学习导图}\par
\vspace{7mm}
\noindent{\sffamily\fontsize{25}{30}\selectfont\bfseries 第 11 天 · 卷积、滤波与池化（宽松排版型）}\par
\vspace{4mm}\noindent\rule{\linewidth}{.85pt}\par\vspace{6mm}
\begin{center}
\begin{tikzpicture}[x=1cm,y=1cm,>=stealth,every node/.style={font=\sffamily}]
\path[use as bounding box] (-8.1,-18.7) rectangle (8.1,.4);
\node[font=\sffamily\Large\bfseries] at (0,-.6) {从像素窗口到特征图};
\node[font=\sffamily\small,text=black!55] at (0,-1.25) {先数有几个窗口，再算每个窗口输出什么};
\node[anchor=north west,font=\sffamily\fontsize{22}{24}\selectfont\bfseries,text=black!50] at (-7.35,-2.35) {01};
\node[anchor=north west,text width=5.8cm,align=left,font=\sffamily\large\bfseries] at (-6.199999999999999,-2.3200000000000003) {像素网格};
\node[anchor=north west,text width=5.9cm,align=left,font=\sffamily\small] at (-6.199999999999999,-3.0700000000000003) {采样决定取哪些位置\\[.25em]量化决定用哪些亮度值};
\draw[black!35,line width=.4pt] (-7.35,-5.050000000000001) -- (-0.2999999999999998,-5.050000000000001);
\node[anchor=north west,font=\sffamily\fontsize{22}{24}\selectfont\bfseries,text=black!50] at (0.85,-2.35) {02};
\node[anchor=north west,text width=5.8cm,align=left,font=\sffamily\large\bfseries] at (2.0,-2.3200000000000003) {窗口位置};
\node[anchor=north west,text width=5.9cm,align=left,font=\sffamily\small] at (2.0,-3.0700000000000003) {填充改变边界，步长决定落点\\[.25em]只数能完整放下核的位置};
\draw[black!35,line width=.4pt] (0.85,-5.050000000000001) -- (7.8999999999999995,-5.050000000000001);
\node[anchor=north west,font=\sffamily\fontsize{22}{24}\selectfont\bfseries,text=black!50] at (0.85,-7.7) {03};
\node[anchor=north west,text width=5.8cm,align=left,font=\sffamily\large\bfseries] at (2.0,-7.67) {权值共享};
\node[anchor=north west,text width=5.9cm,align=left,font=\sffamily\small] at (2.0,-8.42) {每个窗口用同一个核\\[.25em]对应位置相乘，再求和};
\draw[black!35,line width=.4pt] (0.85,-10.4) -- (7.8999999999999995,-10.4);
\node[anchor=north west,font=\sffamily\fontsize{22}{24}\selectfont\bfseries,text=black!50] at (-7.35,-7.7) {04};
\node[anchor=north west,text width=5.8cm,align=left,font=\sffamily\large\bfseries] at (-6.199999999999999,-7.67) {滤波结果};
\node[anchor=north west,text width=5.9cm,align=left,font=\sffamily\small] at (-6.199999999999999,-8.42) {平均核平滑邻域\\[.25em]差分核响应亮度突变};
\draw[black!35,line width=.4pt] (-7.35,-10.4) -- (-0.2999999999999998,-10.4);
\node[anchor=north west,font=\sffamily\fontsize{22}{24}\selectfont\bfseries,text=black!50] at (-7.35,-13.05) {05};
\node[anchor=north west,text width=5.8cm,align=left,font=\sffamily\large\bfseries] at (-6.199999999999999,-13.020000000000001) {激活顺序};
\node[anchor=north west,text width=5.9cm,align=left,font=\sffamily\small] at (-6.199999999999999,-13.770000000000001) {留存卷积结果，再按题意做 ReLU\\[.25em]先置零会改变后面的池化输入};
\draw[black!35,line width=.4pt] (-7.35,-15.75) -- (-0.2999999999999998,-15.75);
\node[anchor=north west,font=\sffamily\fontsize{22}{24}\selectfont\bfseries,text=black!50] at (0.85,-13.05) {06};
\node[anchor=north west,text width=5.8cm,align=left,font=\sffamily\large\bfseries] at (2.0,-13.020000000000001) {池化摘要};
\node[anchor=north west,text width=5.9cm,align=left,font=\sffamily\small] at (2.0,-13.770000000000001) {每个通道独立取最大或平均\\[.25em]窗口大小与步长决定输出尺寸};
\draw[black!35,line width=.4pt] (0.85,-15.75) -- (7.8999999999999995,-15.75);
\draw[->,black!55,line width=.8pt] (-.12,-3.7) -- (.55,-3.7);
\draw[->,black!55,line width=.8pt] (7.85,-5.45) -- (7.85,-7.2);
\draw[->,black!55,line width=.8pt] (.52,-9.05) -- (-.15,-9.05);
\draw[->,black!55,line width=.8pt] (-7.85,-10.8) -- (-7.85,-12.55);
\draw[->,black!55,line width=.8pt] (-.12,-14.4) -- (.55,-14.4);
\node[font=\sffamily\small,text=black!55] at (0,-18.1) {10 分钟热身 · 25 分钟讲解 · 40 分钟练习 · 15 分钟订正};
\end{tikzpicture}
\end{center}
\clearpage

\chaptertitle{第 11 天 · 卷积、滤波与池化}
\textbf{2026 年 11 月 25 日｜K12–K14｜90 分钟}

把卷积核想成一块可以移动的透明模板。模板盖住几格像素，同一位置的数相乘，再把乘积加起来。模板每挪一次，就在特征图上填一个数。接着，池化可以把特征图的一小块浓缩成一个数。今天要独立算出两张特征图和池化结果，并用一句话解释成像或视差成立的条件。

\section{90 分钟安排}
\begin{itemize}
\item 10 分钟 · H1 复做隐藏层梯度（3 分钟），H2 回顾权值共享（7 分钟）。
\item 25 分钟 · 看两个单元的例题，做 A1、B1。需要的前置只有矩阵的行列、对应位置相乘后求和，以及 ReLU 把负数变成 0。
\item 40 分钟 · A2、B2 各 7 分钟，F1 用 6 分钟，M1、M2 各 10 分钟。
\item 15 分钟 · 对照解析，圈出自己最先出错的一步。合上答案重算一个窗口，再口头解释一个视觉条件。
\end{itemize}
H1 来自本课程错题本第 22 条的复做题。其余为新编题，覆盖备考计划中本日的填空和计算考点。

\subsection{举一反三 · 热身}
\PracticeQuestion{H1}{输入 $x=(2,-1)^\top$，两个神经元激活前的梯度已经算出，为 $\delta=(0,3)^\top$。用 $\delta_j x_i$ 写出第 $j$ 个神经元接收第 $i$ 个输入的权重梯度。写出完整的权重梯度矩阵和偏置梯度，并分别标注形状。}{35}
\PracticeQuestion{H2}{一张 $5\times5$ 单通道图，用一个 $2\times2$ 核，步长 1、填充 0、每个核带一个偏置。左上窗口和它右边的窗口，分别读取哪两列？两个窗口共用几项可学习参数？说出「局部连接」和「权值共享」各指什么。}{35}
\clearpage
\section{单元 1 · 卷积核滑一次，特征图填一格}
核只读当前窗口，这是\textbf{局部连接}。同一个核在各个窗口重复使用同一组权重及偏置，这是\textbf{权值共享}。本课用常见 CNN 约定，核按给定方向直接相乘（互相关）。若题目要求数学卷积，先将核旋转 $180^\circ$。本课卷积偏置均为 0，膨胀率为 1。

步长 $S$ 是相邻窗口左上角之间的距离。若四边各补 $P$ 圈零，沿高、宽分别套式，$N$ 是输入边长，$K$ 是核在该方向的边长。

\[O=\left\lfloor\frac{N+2P-K}{S}\right\rfloor+1\qquad(S\ge1,\ N+2P\ge K).\]
分子表示核左上角最多能移动多少格。只数完整窗口，所以向下取整。末尾的 1 对应起始窗口。填充可让边缘参与更多窗口，输出大小仍要代公式判断。

\subsection{完整例题 · 用同一个核算四格}
输入 $X=\begin{pmatrix}1&2&0\\0&1&3\\2&1&0\end{pmatrix}$，核 $K=\begin{pmatrix}1&0\\0&-1\end{pmatrix}$，步长 1、填充 0。求激活前的特征图。

\begin{enumerate}
\item 行列各有 $\lfloor(3-2)/1\rfloor+1=2$ 格，所以先画 $2\times2$。
\item 核只有左上和右下的权重非零。左上窗口是 $\begin{pmatrix}1&2\\0&1\end{pmatrix}$，对应位置乘加后为 $1-1=0$。
\item 向右滑一格算 $2-3=-1$。从原位置向下滑一格算 $0-1=-1$，再向右滑一格算 $1-0=1$。结果为 $Y=\begin{pmatrix}0&-1\\-1&1\end{pmatrix}$。核每次都做「左上减右下」，可以用这一点检查四格。
\end{enumerate}
\subsection{举一反三}
\PracticeQuestion{A1}{沿用例题的核和设置，仅把输入右下角从 0 改为 4。前三个输出已知仍为 $0,-1,-1$。补出右下输出的乘加，说明为何只有这一格改变。}{23}
\PracticeQuestion{A2}{恢复例题的原始 $X$，核不变。把步长改为 2，四边各补 1 圈零。先求输出尺寸，再写四个窗口的左上角坐标和特征图。坐标以补零后矩阵的第 1 行、第 1 列为起点。指出哪些窗口用到了补入的零。}{26}
\clearpage
\section{单元 2 · 池化给每个窗口做一个摘要}
卷积计算每个窗口的加权和。池化对已有数值取最大值或平均值，常用来缩小空间尺寸。多通道时，每张特征图各自池化。给定窗口和步长，最大池化、平均池化本身都没有可学习权重。本课池化不填充，且只处理完整窗口，可用单元 1 的尺寸公式令 $P=0$。

\subsection{完整例题 · 同一张特征图，两种摘要}
取单元 1 的 $Y=\begin{pmatrix}0&-1\\-1&1\end{pmatrix}$，直接池化。窗口 $2\times2$、步长 2。分别求最大和平均池化输出。

\begin{enumerate}
\item 每边长度恰好等于窗口尺寸，只有一个完整窗口。输出大小为 $1\times1$。
\item 最大池化保留窗口里最强的数。最大值是 $\max(0,-1,-1,1)=1$。
\item 平均池化用窗口所有值，结果为 $(0-1-1+1)/4=-1/4$。用范围验算，均值落在最小值 $-1$ 和最大值 1 之间。如果先用 ReLU，窗口内的负值会被改成 0，均值将随之改变。
\end{enumerate}
\subsection{举一反三}
\PracticeQuestion{B1}{例题改为先 ReLU 再平均池化。已知 ReLU 输出为 $\begin{pmatrix}0&0\\0&1\end{pmatrix}$，补出唯一一个均值并标注输出大小。}{24}
\PracticeQuestion{B2}{输入 $Q=\begin{pmatrix}-4&-1&2\\-3&-2&0\\1&4&-5\end{pmatrix}$，直接做 $2\times2$ 最大池化，步长 1。先写尺寸，再求完整输出。额外算出左上窗口的平均值，解释左上最大值为何可以为负。本题窗口会重叠，说明与例题的区别。}{34}
\clearpage
\section{视觉速记 · 像素怎么来，滤波能看出什么}
\begin{itemize}
\item \textbf{成像。} 在相机坐标中，以光心为原点，使用光心前方的虚拟成像平面（距离为焦距 $f$）。点 $(X,Y,Z)$ 在相机前方 $Z>0$ 时，投影为 $(u,v)=(fX/Z,fY/Z)$，这里坐标以主点为原点，轴方向一致。固定 $X,Y$ 时，增大 $Z$ 会使投影更靠近主点。采样选取空间位置，把图像变成网格。量化让像素值落入有限的亮度级别。
\item \textbf{滤波。} 归一化平均核让邻近像素互相平均，能减弱小波动，也会模糊边缘。归一化高斯核给近邻较大权重，通常尺度越大，保留的细节越少。权重和为 0 的差分核在常量窗口输出 0，在有方向性亮度变化的窗口可能产生较大的正值或负值。
\item \textbf{视差。} 水平校正后、焦距相同的平行双目中，左右图的同一点需要正确匹配。若两图的横坐标已各自减去主点横坐标，取 $d=u_L-u_R>0$，深度为 $Z=fB/d$，$B$ 为基线长度。若 $f$ 以像素为单位，$d$ 也要用像素。同一套相机下，近点视差大。重复纹理等情况会让多个窗口同样相似，匹配错误时，代入数值也会得到错误的深度。
\end{itemize}
\section{独立练习}
\textbf{先独立作答，再翻到后面的解析。} 计算题至少写一个窗口完整的乘加，其余窗口要能看出对应的输入位置。

\PracticeQuestion{F1}{两个小填空。（1） 把连续图像变为像素网格叫作哪一步？把亮度取到 256 个级别叫作哪一步？（2） 水平校正、坐标已减主点且匹配正确的一套双目，$f=100$ 像素，$B=0.12$ 米。两个点的视差分别是 12 和 6 像素，各多远？哪个近？}{35}
\clearpage
\PracticeQuestion{M1}{输入 $U=\begin{pmatrix}1&1&0&2\\1&1&0&2\\3&1&2&2\\3&1&2&2\end{pmatrix}$，核 $E=\begin{pmatrix}1&-1\\1&-1\end{pmatrix}$。步长 2、填充 0。写出窗口位置和激活前的特征图，再经 ReLU 和 $2\times2$、步长 2 的最大池化。标注各步大小。解释右上和左下窗口的卷积值为何符号相反。}{64}
\PracticeQuestion{M2}{输入 $V=\begin{pmatrix}2&2&6\\2&2&6\\2&2&6\end{pmatrix}$。给定平均核 $G=\frac14\begin{pmatrix}1&1\\1&1\end{pmatrix}$ 和 M1 的 $E$。任务是找亮度从左往右突变的位置。两种核都用步长 1、填充 0，分别计算两张特征图，再说明应选哪一个。若把 $V$ 改成全 2，两个核各输出什么？}{64}
\clearpage
\section{答案与解析 · 热身和单元题}
\textbf{H1.} 第 1 个神经元给出 $0(2,-1)=(0,0)$，第 2 个给出 $3(2,-1)=(6,-3)$。合成 $\frac{\partial L}{\partial W}=\begin{pmatrix}0&0\\6&-3\end{pmatrix}$，形状 $2\times2$。每个偏置只直接加到自己的激活前数值，故偏置梯度是 $\delta=(0,3)^\top$，形状 $2\times1$。

\textbf{H2.} 左上窗口读第 1、2 列，右邻窗口读第 2、3 列，两者都读第 1、2 行。每个输出只看当前 4 格，是局部连接。两个窗口共用同一个核的 4 个权重和 1 个偏置，共 5 项可学习参数，是权值共享。

\textbf{A1.} 右下输入只属于右下的完整窗口，所以只有一个输出变化。右下窗口为 $\begin{pmatrix}1&3\\1&4\end{pmatrix}$，乘加为 $1(1)+3(0)+1(0)+4(-1)=-3$。

\textbf{A2.} 补零后是 $5\times5$，输出每边 $\lfloor(5-2)/2\rfloor+1=2$。窗口左上角依次为 $(1,1),(1,3),(3,1),(3,3)$。仍按左上减右下，四格依次为 $0-1=-1$、$0-0=0$、$0-2=-2$、$1-0=1$，输出 $\begin{pmatrix}-1&0\\-2&1\end{pmatrix}$。前三个窗口用到补零圈，最后一个只读原矩阵。最后一个减去的 0 来自原矩阵。坐标不会走到第 5 行或第 5 列，因为那样放不下完整窗口。

\textbf{B1.} 先经过 ReLU，负值都变 0，四数和为 1，均值为 $1/4$，输出 $1\times1$。

\textbf{B2.} 每边 $\lfloor(3-2)/1\rfloor+1=2$。各窗口按行读取的数据分别是 $( -4,-1,-3,-2)$、$(-1,2,-2,0)$、$(-3,-2,1,4)$ 和 $(-2,0,4,-5)$，逐窗取最大值得 $\begin{pmatrix}-1&2\\4&4\end{pmatrix}$。左上均值为 $(-4-1-3-2)/4=-2.5$，四个输入都负，最大值也负。步长从 2 改为 1，使相邻的 $2\times2$ 窗口有重叠。

\clearpage
\subsection{混合题答案}
\textbf{F1.} （1） 采样、量化。（2） $100(0.12)/12=1$ 米，$100(0.12)/6=2$ 米。相机相同且匹配、坐标约定成立，所以视差 12 像素的点更近。

\textbf{M1.} 每边 $\lfloor(4-2)/2\rfloor+1=2$。左上角按行分别在 $(1,1),(1,3),(3,1),(3,3)$。核做「左列和减右列和」，按行得到 $(1+1)-(1+1)=0$、$(0+0)-(2+2)=-4$、$(3+3)-(1+1)=4$、$(2+2)-(2+2)=0$。右上是右列更亮，左下是左列更亮，所以异号。卷积结果 $\begin{pmatrix}0&-4\\4&0\end{pmatrix}$ 经 ReLU 得 $\begin{pmatrix}0&0\\4&0\end{pmatrix}$，两者都为 $2\times2$。全图正好是一个池化窗口，最大值 4，输出 $1\times1$。

\textbf{M2.} 两个核都得到 $2\times2$ 输出。两行的窗口结构相同。左窗四个数都是 2，右窗两列分别为 $(2,2)^\top$、$(6,6)^\top$。均值分别为 $(2+2+2+2)/4=2$ 和 $(2+6+2+6)/4=4$，输出 $\begin{pmatrix}2&4\\2&4\end{pmatrix}$。差分分别为 $4-4=0$ 和 $4-12=-8$，输出 $\begin{pmatrix}0&-8\\0&-8\end{pmatrix}$。选差分核，大幅响应直接标出有突变的窗口，负号反映右边更亮。输入全为 2 时，平均核四格都输出 2，差分核四格都输出 0，因为平均核权重和是 1，差分核权重和是 0。

\section{完成标准与反馈}
合上解析后，检查自己能否先定尺寸再逐窗计算，独立写出 M1 和 M2 中的特征图以及 M1 的池化结果，并用一句话说清一个成像或视差条件。复习时先画窗口并确定落点，再对应位置乘加，最后按题目指定的次序做 ReLU 和池化。

可以发 H1、M1、M2 的计算过程或手写照片。附上实际用时、哪一步靠提示完成，便于决定下一课热身要复做哪一步。

\end{document}
